% Source-faithful assembly: Mumford, Abelian Varieties, Appendix I,
% Horizon transcription pages 0252--0261 (Theorem 1, Steps I--VII).

\section{Tate's theorem over finite fields}
\label{mumford-appendix-I-source}
\dcref{M:appI:I}

\begin{theorem}[Tate]
\label{mumford-thm-tate-hom-finite-field}
\label{mumford-thm-tate-hom-isomorphism}
Let $X$ and $Y$ be abelian varieties defined over a finite field $k_0$. Then
the homomorphism
\[
 \mathbb Z_l\otimes_{\mathbb Z}{\rm Hom}(X,Y)\longrightarrow
 {\rm Hom}_{\mathbb Z_l}(T_l(X),T_l(Y))^{(G)}
\]
is an isomorphism, where $G={\rm Gal}(k/k_0)$ and the superscript denotes
elements fixed by some neighborhood of the identity.
\source{mumford-abelian-varieties:page-0252}
\uses{mumford-thm-hom-tate-injective}
\uses{mumford-appI-step-I,mumford-appI-step-II,mumford-appI-step-III,
 mumford-appI-step-IV,mumford-appI-step-V,mumford-appI-step-VI,
 mumford-appI-step-VII}
\end{theorem}

\begin{proposition}[Step I]
\label{mumford-appI-step-I}
It suffices to show that
\[
 \mathbb Q_l\otimes_{\mathbb Z_l}{\rm Hom}(X,Y)
 =\mathbb Q_l\otimes_{\mathbb Q}{\rm Hom}^{0}(X,Y)
 \longrightarrow {\rm Hom}_{\mathbb Q_l}(V_l(X),V_l(Y))^{(G)} \tag{3}
\]
is bijective, where $V_l(X)=\mathbb Q_l\otimes_{\mathbb Z_l}T_l(X)$.
\source{mumford-abelian-varieties:page-0252}
\end{proposition}
\begin{proof}
If (3) is bijective, the image $M$ of (1) is a $\mathbb Z_l$-submodule of
maximal rank in ${\rm Hom}_{\mathbb Z_l}(T_l(X),T_l(Y))^{(G)}$. We claim that
${\rm Hom}_{\mathbb Z_l}(T_l(X),T_l(Y))/M$ has no torsion. Suppose that
$\varphi$ is a homomorphism and $l\varphi\in M$. Choose
$\psi_n\in {\rm Hom}(X,Y)$ such that $T_l(\psi_n)\to l\varphi$.
Then $T_l(\psi_n)(T_l(X))\subset lT_l(Y)$, so $\psi_n$ vanishes on $X_l$.
Consequently $\psi_n$ factors as
\[
 \psi_n=\chi_n\circ l_Y=l\chi_n.
\]
The $\chi_n$ converge to an element $\chi$ of
$\mathbb Z_l\otimes_{\mathbb Z}{\rm Hom}(X,Y)$ with $T_l(\chi)=\varphi$.
Thus $\varphi\in M$, proving the claim. Therefore $M$ is the whole invariant
module, as required.

\end{proof}


\begin{proposition}[Step II]
\label{mumford-appI-step-II}
It suffices to show, for every abelian variety $X$ over a finite field, that
\[
 \mathbb Q_l\otimes_{\mathbb Q}{\rm End}^{0}(X)
 \xrightarrow{\lambda}{\rm End}_{\mathbb Q_l}(V_l(X))^{(G)} \tag{4}
\]
is an isomorphism.
\source{mumford-abelian-varieties:page-0253}
\end{proposition}
\begin{proof}
Assume (4) is an isomorphism for $X\times Y$. The direct-sum decompositions
as $G$-modules are
\[
\begin{aligned}
 \mathbb Q_l\otimes {\rm End}^{0}(X\times Y)&\simeq
 (\mathbb Q_l\otimes {\rm End}^{0}X)\oplus(\mathbb Q_l\otimes {\rm End}^{0}Y)\\
 &\quad\oplus(\mathbb Q_l\otimes {\rm Hom}^{0}(X,Y))\oplus
 (\mathbb Q_l\otimes {\rm Hom}^{0}(Y,X)),
\end{aligned}
\]
and
\[
 {\rm End}(V_l(X\times Y))\simeq {\rm End}(V_l(X))\oplus {\rm End}(V_l(Y))
 \oplus {\rm Hom}(V_l(X),V_l(Y))\oplus {\rm Hom}(V_l(Y),V_l(X)).
\]
The resulting diagram is commutative and its first vertical arrow is an
isomorphism; hence (3) is an isomorphism.

\end{proof}


\begin{proposition}[Step III]
\label{mumford-appI-step-III}
It suffices to show that $\lambda$ is an isomorphism for one $l$ and that
\[\dim_{\mathbb Q_l}{\rm End}(V_l(X))^{(G)}\]
is independent of $l\ne {\rm char}\,k$.
\source{mumford-abelian-varieties:page-0253}
\end{proposition}
\begin{proof}
The dimension of the left member of (4) is independent of $l$, and $\lambda$
is always injective.

\end{proof}


\begin{proposition}[Step IV]
\label{mumford-appI-step-IV}
Let $E$ be the image of $\lambda$, and let $F$ be the intersection over all
neighborhoods of $e$ in $G$ of the subalgebras generated in
${\rm End}_{\mathbb Q_l}(V_l)$ by these neighborhoods. To establish (4), it
suffices to show that $F$ is the commutant of $E$ in ${\rm End}(V_l)$.
\source{mumford-abelian-varieties:page-0253}
\end{proposition}
\begin{proof}
Since $E$ is semisimple, von Neumann's density theorem then implies that $E$
is the commutant of $F$. Since ${\rm End}_{\mathbb Q_l}(V_l)$ is finite
dimensional over $\mathbb Q_l$, $F$ equals the subalgebra generated by some
neighborhood, and all smaller neighborhoods generate the same subalgebra.
Hence its commutant is precisely ${\rm End}_{\mathbb Q_l}(V_l)^{(G)}$.

\end{proof}


\begin{convention}[Hypothesis $(k_0,X,l)$]
\label{mumford-appI-hypothesis}
Let $X$ be an abelian variety over $k_0$ and $l$ a prime. For every integer
$d\geq1$, there are, up to isomorphism, only finitely many abelian varieties
$Y$ over $k_0$ such that (a) $Y$ has an ample line bundle $L$ defined over
$k_0$ with $\chi(L)=d$, and (b) there is a $k_0$-isogeny $Y\to X$ of $l$-power
degree.
\source{mumford-abelian-varieties:page-0254}
\end{convention}

\begin{lemma}[Finite-field finiteness]
\label{mumford-lem-appI-finite-field-finiteness}
If $k_0$ is finite, $d>0$, and $g>0$, only finitely many isomorphism classes
of $g$-dimensional abelian varieties $Y$ over $k_0$ carry an ample line bundle
$L$ defined over $k_0$ with $\chi(L)=d$.
\source{mumford-abelian-varieties:page-0254}
\end{lemma}
\begin{proof}
The line bundle $L^3$ gives a projective embedding of $Y$ in projective space
of dimension $3^g\chi(L)-1=3^gd-1$. Its degree is the $g$-fold self
intersection number of $L^3$, namely
\[
 3^g(L)^g=3^g g!d.
\]
Thus it gives a $k_0$-rational point of the Chow variety of cycles of
dimension $g$ and degree $3^gg!d$ in $\mathbb P^{3^gd-1}$. A finite field has
only finitely many rational points on this Chow variety.

\end{proof}


\begin{lemma}[Quotient descent]
\label{mumford-lem-appI-quotient-descent}
Let $X$ be quasi-projective over $k_0$ and $\Pi$ a finite group of
automorphisms such that (i) the automorphisms of $\Pi$ are defined over a
separably algebraic extension of $k_0$, and (ii) $\Pi^\sigma=\Pi$ for every
automorphism $\sigma$ of $\bar k_0$ over $k_0$. Then $X/\Pi$ and
$X\to X/\Pi$ are defined over $k_0$.
\source{mumford-abelian-varieties:page-0255}
\end{lemma}
\begin{proof}
Choose a Galois extension $k_1/k_0$ over which $\Pi$ is defined. Replacing an
affine open $U$ by $\bigcap_{\lambda\in\Pi}\lambda U$, assume
$X={\rm Spec}\,A$ with $A=A_0\otimes_{k_0}k$. Put
$A_1=A_0\otimes_{k_0}k_1$. Then $A_1^\Pi\otimes_{k_1}k=A^\Pi$.
The group $G={\rm Gal}(k_1/k_0)$ stabilizes $A_1^\Pi$. If $B$ is its
$k_0$-invariant subalgebra and $\{\theta_i\}$ is a basis of $k_1/k_0$, then
for $\sum a_i\otimes\theta_i\in A_1^\Pi$, applying every element of $G$ and
using $\det(\lambda(\theta_i))_{\lambda,i}\ne0$ gives
$a_i\otimes1\in A_1^\Pi\cap A_0\subset B$. Thus
 $A_1^\Pi=B\otimes_{k_0}k_1$, proving descent.
\source{mumford-abelian-varieties:page-0256}

\end{proof}


\begin{lemma}[Galois stability of ${\rm Pic}^0$]
\label{mumford-lem-appI-piczero-galois}
If $L\in{\rm Pic}^0Y$ is defined over $\bar k_0$ and $\sigma\in
{\rm Gal}(\bar k_0/k_0)$, then $\sigma(L)\in{\rm Pic}^0Y$.
\source{mumford-abelian-varieties:page-0257}
\end{lemma}
\begin{proof}
Let $M$ be ample over $k_0$ and put
$N=n_1^*M\otimes p_1^*M^{-1}\otimes p_2^*M^{-1}$ on $Y\times Y$.
Choose $y$ with $N|_{\{y\}\times Y}=L$. Then
$N|_{\{\sigma y\}\times Y}=\sigma(L)$, which is in ${\rm Pic}^0Y$.

\end{proof}


\begin{lemma}[Descent of invariant line bundles]
\label{mumford-lem-appI-line-bundle-descent}
Let $Y$ be complete over $k_0$ with a $k_0$-point, let $k_1/k_0$ be Galois,
and let $L$ on $\overline Y$ be defined over $k_1$ with
$\sigma(L)\simeq L$ for every $\sigma\in{\rm Gal}(k_1/k_0)$. Then $L$ is
defined over $k_0$.
\source{mumford-abelian-varieties:page-0257}
\end{lemma}
\begin{proof}
Choose an affine cover $\{U_i\}$ of $Y_0$ over which $L$ is free. If
$\alpha_{ij}$ is the defining cocycle, invariance gives units $\beta_{i,\sigma}$
with
\[
 \frac{\sigma(\alpha_{ij})}{\alpha_{ij}}=\frac{\beta_{i,\sigma}}{\beta_{j,\sigma}}.
\]
Normalize at a rational point so $\beta_{i_0,\sigma}(y_1)=1$. Comparing the
two expressions for $\sigma\tau(\alpha_{ij})/\alpha_{ij}$ shows
\[
 \beta_{i,\sigma\tau}\sigma(\beta_{i,\tau})^{-1}\beta_{i,\sigma}^{-1}
 =\beta_{j,\sigma\tau}\sigma(\beta_{j,\tau})^{-1}\beta_{j,\sigma}^{-1};
\]
completeness and evaluation at $y_1$ make this equal to $1$. The local-unit
cohomology $H^1(G,(A\otimes_{k_0}k_1)^*)=1$ therefore gives local units
$\gamma_{i,\alpha}$ with
$\beta_{i,\sigma}=\sigma(\gamma_{i,\alpha})/\gamma_{i,\alpha}$.
The modified cocycle
$\alpha_{ij}\gamma_{j,\beta}/\gamma_{i,\alpha}$ is Galois invariant and
therefore comes from $Y_0$.

For completeness, let $B=A\otimes_{k_0}k_1$ for a local ring $A$ of $Y_0$,
and $R$ its quotient field. The exact sequence is
\[
 H^0(G,R^*)\to H^0(G,R^*/B^*)\to H^1(G,B^*)\to H^1(G,R^*)=\{1\}.
\]
If $f\in R^*$ represents an invariant class, write $f=f_0/b$ and assume
$f\in B$. The ideal $Bf$ is Galois invariant, hence $B\mathfrak a$ for an
ideal $\mathfrak a$ of $A$. Since $Bf\simeq k_1\otimes_{k_0}\mathfrak a$,
$\mathfrak a$ is projective and principal; consequently $f=gu$ with
 $g\in A$ invariant and $u\in B^*$. This proves the local assertion and the
 lemma.
\source{mumford-abelian-varieties:page-0258}
\source{mumford-abelian-varieties:page-0259}

\end{proof}


\begin{proposition}[Step V]
\label{mumford-appI-step-V}
Suppose $X$ and an ample line bundle $L$ on $X$ are defined over $k_0$,
$l\ne {\rm char}\,k_0$, and Hypothesis $(k_0,X,l)$ holds. If
$W\subset V_l(X)$ is $G$-stable and maximal isotropic for $E^L$, then there
is an element $u\in F$ with $u(V_l(X))=W$.
\source{mumford-abelian-varieties:page-0254}
\end{proposition}
\begin{proof}
Set $T=T_l(X)$ and $V=V_l(X)$. For $n>0$, put
\[
 T_n=(T\cap W)+l^nT.
\]
If $\psi_n:T_n(X)\to X_n$ is natural, set $K_n=\psi_n(T_n)$,
$Y_n=X/K_n$, and let $\pi_n:X\to Y_n$ be natural. Then
\[
 (l^n)_X\text{ factors as }X\xrightarrow{\pi_n}Y_n\xrightarrow{\lambda_n}X.
\]
Since $\pi_n\lambda_n\pi_n=l^n\pi_n$,
\[
 \lambda_n((Y_n)_{l^n})=\ker\pi_n=K_n.
\]
For $m\ge n$,
$\lambda_n((Y_n)_{l^m})\supset(\lambda_n\pi_n)(X_{l^m})=X_{l^{m-n}}$;
hence
\[
 T_l(\lambda_n)(T_l(Y_n))\supset l^nT_l(X),\qquad
 T_l(\lambda_n)(T_l(Y_n))=T_n.
\]
The points of $X_{l^n}$ are separably algebraic over $k_0$, because
$(l^n)_{X_0}$ is etale. Since $W$ and $T_n$ are $G$-stable, so is $K_n$.

By the quotient descent lemma, $Y_n$ and $\pi_n$ are defined over $k_0$.
The addition and inverse on $Y_n$, defined over a separably algebraic
extension and Galois invariant, descend to $k_0$, so $Y_n$ is an abelian
variety. Since $(l^n)_X=\lambda_n\pi_n$ and $\pi_n$ are defined over $k_0$,
$\lambda_n$ is defined over $k_0$.

Let $d=\chi(L)$. For $x,y\in T_l(Y_n)$,
\[
 e_l(x,\phi_{n_X^*(L)}(y))=E^{n_X^*(L)}(x,y)
 =E^L(\lambda_n(x),\lambda_n(y))\in e^L(T_n,T_n)\subset l^nM.
\]
Non-degeneracy of $e_l$ gives $\phi_{n_X^*(L)}=l^n\psi$ for some
$\psi:\widehat Y_n\to Y_n$. Hence $\psi=\varphi_{L_n}$ for a line bundle
$L_n$; after replacing by a power we may assume $L_n$ is symmetric and
$L_n^{p^N}$ is defined over a Galois extension $k_1/k_0$.

Write $L_1\equiv L_2$ when $L_1\otimes L_2^{-1}\in{\rm Pic}^0Y$. Since
$NS(Y_n)$ is torsion-free, $\sigma(L_n^{p^N})\equiv L_n^{p^N}$; symmetry
then shows their difference has order $2$. Thus $M_n=L_n^{2p^N}$ satisfies
$\sigma(M_n)\equiv M_n$ and
\[
 M_n^{\,n}\simeq\lambda_n^*(L_n^{2p^N})\simeq L_n^{2p^Nn^2}.
\]

The descent lemma makes $M_n$ defined over $k_0$. Choose $a,b$ with
$2ap^N+bl^n=2$ and set
\[
 N_n=M_n^a\otimes\lambda_n^*(L_n)^b.
\]
Then $N_n$ is ample over $k_0$, $N_n\simeq L_n^2$, and
\[
 \chi(N_n)=2^g\chi(L_n)=2^g l^{-ng}\chi(\lambda_n^*L)
 =2^g l^{-ng}\deg(\lambda_n)\chi(L)=2^gd.
\]
The Poincare-bundle construction gives the same conclusion and
$\chi(N_n)^2=2^{2g}\chi(L)^2$.

By Hypothesis there is an infinite set of $n$ giving isomorphic $Y_n$; choose
the smallest $n$ and isomorphisms $v_i:Y_n\to Y_i$. Put
$u_i=\lambda_i v_i\lambda_n^{-1}$. Their images satisfy
$u_i'(T_n)\subset T_i\subset T_n$. Compactness of
${\rm End}_{\mathbb Z_l}(T_n)$ gives a convergent subsequence with limit
$u'\in E$. Since the $T_i$ decrease,
\[
 u'(T_n)=\bigcap_i u_i'(T_n)=\bigcap_iT_i=T_n\cap W.
\]
Thus $u'(V)=W$.
\end{proof}

\begin{proposition}[Step VI]
\label{mumford-appI-step-VI}
Suppose Hypothesis $(k_1,X,l)$ holds for every finite algebraic extension
$k_1/k_0$, and $F$ is a product of copies of $\mathbb Q_l$. Then (4) is an
isomorphism.
\source{mumford-abelian-varieties:page-0260}
\uses{mumford-appI-step-V,mumford-appI-hypothesis}
\end{proposition}
\begin{proof}
Extend $k_0$ so all elements of ${\rm End}\,X$ are defined and
${\rm Gal}(\bar k_0/k_0)$ generates $F$. Let $D$ be the commutant of $E$;
$D\supset F$. We show by downward induction on $\dim W$ that every
$F$-stable isotropic $W$ for $E^L$ is $D$-stable. If $\dim W=g$, Step V gives
$u\in E$ with $u(V)=W$, so $DW=DuV=uDV=uV=W$. If $\dim W=r<g$, then
$W^\perp$ is $F$-stable and
\[
 \dim W^\perp-\dim W=2g-2r\ge2.
\]
Choose $F$-stable lines $L_1,L_2$ with
$W\oplus L_1\oplus L_2$ direct. By induction, $W\oplus L_i$ is $D$-stable,
so their intersection $W$ is $D$-stable. Every $F$-eigenvector is therefore
a $D$-eigenvector, whence $D\subset F$; thus $F=D$.

\end{proof}


\begin{proposition}[Step VII: end of proof]
\label{mumford-appI-step-VII}
After a finite extension assume every element of ${\rm End}\,X$ is defined.
For Frobenius $\pi$, the algebra $\mathbb Q[\pi]$ is commutative semisimple.
There are infinitely many primes $l$ for which
$\mathbb Q_l\otimes_{\mathbb Q}\mathbb Q[\pi]$ is a product of copies of
$\mathbb Q_l$.
\source{mumford-abelian-varieties:page-0261}
\uses{mumford-appI-step-VI}
\end{proposition}
\begin{proof}
Write $\mathbb Q[\pi]=K_1\times\cdots\times K_r$ and choose a finite Galois
extension $K$ containing all $K_i$. It is enough that
$K\otimes_{\mathbb Q}\mathbb Q_l$ have a factor $\mathbb Q_l$. The Galois group
acts transitively on the simple factors: otherwise a proper sum of factors
would give an idempotent fixed by the Galois group, whereas only diagonal
elements $(\alpha,\ldots,\alpha)$ with $\alpha\in\mathbb Q_l$ are fixed.
Choose an algebraic integer $\alpha$ generating $K$, with monic minimal
polynomial $F(X)\in\mathbb Z[X]$. For primes $l$ not dividing its discriminant,
an integer root $n$ of $F$ modulo $l$ is simple and Hensel's lemma lifts it to
a root in $\mathbb Z_l$. Thus it remains to prove the following lemma.

\end{proof}


\begin{lemma}[Infinitely many prime roots]
\label{mumford-lem-appI-prime-root}
For every non-constant $F(X)\in\mathbb Z[X]$, there are infinitely many primes
$l$ for which $F(X)\equiv0\pmod l$ has a solution in $\mathbb Z$.
\source{mumford-abelian-varieties:page-0261}
\uses{mumford-appI-step-VII}
\end{lemma}
